% =====================================================================
\section{Results}
\label{sec:results}

% Results by research question (Framework section). Tables from scripts/make_rq_tables.py; paired counts from
% results/analysis_tuned/paired.csv (item bootstrap within each window, 90% intervals).
\newcommand{\rqAvgBestCoarse}{8}
\newcommand{\rqAvgBestItem}{9}
\newcommand{\rqAvgBestBlock}{5}
\newcommand{\rqAvgBestCoarseWord}{eight}
\newcommand{\rqAvgBestItemWord}{nine}
\newcommand{\rqAvgBestBlockWord}{five}
\newcommand{\rqCells}{27}
\newcommand{\rqAvgBestCells}{22}
\newcommand{\rqAvgWorstCells}{0}
\newcommand{\rqWmWorstCells}{16}
\newcommand{\rqLgbWorstCells}{11}
\newcommand{\rqAvgSecondBlock}{4}
\newcommand{\rqBlockMfiveWm}{0.598}
\newcommand{\rqBlockMfiveLgb}{0.953}
\newcommand{\rqBlockMfiveAvg}{0.589}
\newcommand{\rqBlockFavWm}{0.482}
\newcommand{\rqBlockFavLgb}{0.565}
\newcommand{\rqBlockFavAvg}{0.412}
\newcommand{\rqBlockSdrWm}{0.576}
\newcommand{\rqBlockSdrLgb}{0.632}
\newcommand{\rqBlockSdrAvg}{0.470}
\newcommand{\rqPAvgWmBetter}{128}
\newcommand{\rqPAvgWmWorse}{27}
\newcommand{\rqPAvgWmN}{312}
\newcommand{\rqPAvgWmCoarsebetter}{68}
\newcommand{\rqPAvgWmCoarseworse}{19}
\newcommand{\rqPAvgWmCoarsen}{212}
\newcommand{\rqPAvgWmItembetter}{60}
\newcommand{\rqPAvgWmItemworse}{8}
\newcommand{\rqPAvgWmItemn}{100}
\newcommand{\rqPAvgLgbBetter}{124}
\newcommand{\rqPAvgLgbWorse}{21}
\newcommand{\rqPAvgLgbN}{312}
\newcommand{\rqPAvgLgbCoarsebetter}{77}
\newcommand{\rqPAvgLgbCoarseworse}{17}
\newcommand{\rqPAvgLgbCoarsen}{212}
\newcommand{\rqPAvgLgbItembetter}{47}
\newcommand{\rqPAvgLgbItemworse}{4}
\newcommand{\rqPAvgLgbItemn}{100}
\newcommand{\rqPWmLgbBetter}{64}
\newcommand{\rqPWmLgbWorse}{75}
\newcommand{\rqPWmLgbN}{312}
\newcommand{\rqPWmLgbCoarsebetter}{40}
\newcommand{\rqPWmLgbCoarseworse}{45}
\newcommand{\rqPWmLgbCoarsen}{212}
\newcommand{\rqPWmLgbItembetter}{24}
\newcommand{\rqPWmLgbItemworse}{30}
\newcommand{\rqPWmLgbItemn}{100}
\newcommand{\rqAvgItemWorseMax}{8}
\newcommand{\rqMcsWorldModel}{35}
\newcommand{\rqMcsLightgbm}{37}
\newcommand{\rqMcsAverage}{51}
\newcommand{\rqMcsEnsemble}{41}
\newcommand{\rqMcsTimesfm}{12}
\newcommand{\rqMcsCells}{60}
\newcommand{\rqFastShareDomLo}{11}
\newcommand{\rqFastShareDomHi}{21}
\newcommand{\rqFastShareRestLo}{37}
\newcommand{\rqFastShareRestHi}{42}
\newcommand{\rqShortShareDomLo}{11}
\newcommand{\rqShortShareDomHi}{21}
\newcommand{\rqShortShareRestMax}{2}
\newcommand{\rqAvgBestFast}{eight}
\newcommand{\rqAvgBestSlow}{eight}
\newcommand{\rqShortWmBeatsLgb}{eight}
\newcommand{\rqShortFavWm}{0.539}
\newcommand{\rqShortFavLgb}{0.667}
\newcommand{\rqShortMfiveWm}{0.866}
\newcommand{\rqShortMfiveLgb}{1.024}
\newcommand{\rqNineWord}{nine}
\newcommand{\rqSnaWm}{1.009}
\newcommand{\rqSnaLgb}{0.865}
\newcommand{\rqSnaAvg}{0.884}
\newcommand{\rqRuleLevel}{eight}
\newcommand{\rqRuleSna}{0.970}
\newcommand{\rqCostWmItemLo}{0.008}
\newcommand{\rqCostWmItemHi}{0.040}
\newcommand{\rqCostLgbItemLo}{0.058}
\newcommand{\rqCostLgbItemHi}{0.311}
\newcommand{\rqFaWmItemLo}{89}
\newcommand{\rqFaWmItemHi}{104}
\newcommand{\rqFaLgbItemLo}{66}
\newcommand{\rqFaLgbItemHi}{93}
\newcommand{\rqFaAvgItemLo}{78}
\newcommand{\rqFaAvgItemHi}{96}
\newcommand{\rqWmFaItemOff}{11}
\newcommand{\rqKeptSets}{five}
\newcommand{\rqItemWmBeatsAvg}{five}
\newcommand{\rqStoreWmBeatsLgb}{four}
\newcommand{\rqStoreAvgBeatsWm}{one}
\newcommand{\rqStoreWmBeatsAvg}{one}
\newcommand{\rqStoreLevel}{three}
\newcommand{\rqLgbFaItemLoFrac}{two thirds}
\newcommand{\rqCarriedWm}{0.447}
\newcommand{\rqOwnWm}{0.428}
\newcommand{\rqCarriedLgb}{0.450}
\newcommand{\rqOwnLgb}{0.448}
\newcommand{\rqCarriedWmIxs}{0.002}
\newcommand{\rqCarriedLgbIxs}{0.001}
\newcommand{\rqPoolGainMax}{0.007}
\newcommand{\rqPoolGainWm}{0.005}
\newcommand{\rqPoolGainWmSet}{analgesics}
\newcommand{\rqPoolGainLgb}{0.007}
\newcommand{\rqPoolGainLgbSet}{frozen entrees}
\newcommand{\rqPoolLossWm}{0.141}
\newcommand{\rqPoolLossWmSet}{soft drinks}
\newcommand{\rqPoolLossLgb}{0.235}
\newcommand{\rqPoolLossLgbSet}{analgesics}
\newcommand{\rqPoolGainMaxTwo}{0.01}


The results follow the four research questions. Every number comes from the
tuned world model and the tuned LightGBM, each set by the same search on each
data set (Section~\ref{sec:loop}), scored on the paper's windows. The average is
the equal-weight mean of the two forecasts.

\subsection{RQ1: grain}
\label{sec:rq-grain}

\textbf{Use the average at every grain.} Across the nine data sets and the
three views of Table~\ref{tab:rq-grain} (store grain and above by period, item
grain by period, and the whole window as one block), the average has the lowest
error in 22 of 27 cases and is never the least accurate of the three. The world
model is the least accurate in 16 cases and LightGBM in 11.

The paired tests give the same answer. They compare two forecasts level by level
and window by window, and count a difference only when it holds after
resampling items. Against the world model, the average is clearly more accurate
in 128 of 312 comparisons and clearly less accurate in 27; against LightGBM, 124
and 21. At item grain it is clearly less accurate than either model in at most 8
of 100 comparisons. Across data sets and levels it is in the 90\% model
confidence set in \rqMcsAverage{} of \rqMcsCells{} cases, against
\rqMcsLightgbm{} for LightGBM and \rqMcsWorldModel{} for the world model.

\textbf{Grain does not decide between the two models on their own.} The world
model is clearly more accurate than LightGBM in 40 comparisons at store grain
and above and clearly less accurate in 45; at item grain the counts are 24 and
30. Neither is ahead at either grain across the data sets, so a retailer that
picks one model for a grain is picking by data set, not by grain. The average
removes that choice.

\textbf{On the whole promotion block, the average is again the safe choice.} It
is lowest on five of nine data sets and second on the other four. LightGBM is
lowest on analgesics, bottled juice and snacks, but drifts in total where the
world model does not: its block error on M5 is 0.953, against 0.598 for the world
model and 0.589 for the average, and the same pattern shows on Favorita and soft
drinks.

\begin{table*}[t]
\centering\footnotesize\setlength{\tabcolsep}{3pt}
\caption{RQ1, grain. WRMSSE of the tuned world model (WM), the tuned LightGBM (LGB) and their equal-weight average
(Avg), averaged over the levels at store grain and above, over the item levels, and over every level for the whole
forecast window summed as one block; mean over the paper's windows. Lowest of the three in bold.}
\label{tab:rq-grain}
\resizebox{\textwidth}{!}{\begin{tabular}{l ccc ccc ccc}
\toprule
 & \multicolumn{3}{c}{Store and above, by period} & \multicolumn{3}{c}{Item levels, by period} & \multicolumn{3}{c}{Whole block, all levels}\\
\cmidrule(lr){2-4}\cmidrule(lr){5-7}\cmidrule(lr){8-10}
Data set & WM & LGB & Avg & WM & LGB & Avg & WM & LGB & Avg\\
\midrule
Analgesics & 0.688 & 0.698 & \textbf{0.678} & 0.710 & 0.705 & \textbf{0.695} & 0.523 & \textbf{0.458} & 0.472 \\
Bottled juice & 0.569 & 0.566 & \textbf{0.539} & 0.581 & 0.596 & \textbf{0.563} & 0.545 & \textbf{0.404} & 0.445 \\
Cheese & 0.381 & 0.471 & \textbf{0.377} & 0.623 & 0.594 & \textbf{0.578} & 0.445 & 0.425 & \textbf{0.387} \\
Frozen entrees & 0.339 & 0.327 & \textbf{0.320} & 0.375 & 0.350 & \textbf{0.344} & 0.353 & 0.330 & \textbf{0.318} \\
Soft drinks & 0.822 & 0.805 & \textbf{0.754} & 0.589 & 0.555 & \textbf{0.547} & 0.576 & 0.632 & \textbf{0.470} \\
Snacks & 1.241 & \textbf{1.040} & 1.071 & 0.545 & 0.516 & \textbf{0.508} & 0.893 & \textbf{0.633} & 0.688 \\
Frat & 0.471 & 0.479 & \textbf{0.449} & 0.528 & 0.524 & \textbf{0.508} & \textbf{0.345} & 0.406 & 0.364 \\
M5 & 0.660 & 0.634 & \textbf{0.591} & 0.859 & 0.886 & \textbf{0.856} & 0.598 & 0.953 & \textbf{0.589} \\
Favorita & 0.314 & 0.325 & \textbf{0.278} & 0.655 & 0.696 & \textbf{0.653} & 0.482 & 0.565 & \textbf{0.412} \\
\bottomrule
\end{tabular}}
\end{table*}


\subsection{RQ2: nature of the line}
\label{sec:rq-lines}

\textbf{Sales are concentrated, and the average holds up on both the fast lines
and the slow tail; where history is short, the world model is ahead of
LightGBM.} On the six Dominick's categories, 11 to 21\% of the items make 80\% of
unit sales over the year before the forecast; on M5, Frat and Favorita the head
is flatter, at 37 to 42\% (Table~\ref{tab:rq-lines}). On the fast lines the
average has the lowest error on eight of nine data sets, and on the slow lines
also on eight; on snacks LightGBM is lowest on the fast lines, and on cheese the
world model is lowest on the slow lines.

On series whose first sale falls less than 26 weeks before the forecast, the
world model's error is lower than LightGBM's on eight of nine data sets, Frat
being the exception: for example 0.539 against 0.667 on Favorita and 0.866
against 1.024 on M5. These series are 11 to 21\% of the series on the Dominick's
categories, where items enter and leave the range over the years, and under 2\%
on M5, Frat and Favorita.

\begin{table*}[t]
\centering\footnotesize\setlength{\tabcolsep}{3pt}
\caption{RQ2, nature of the line. Share of items that make 80\% of unit sales over the year before the forecast, and
WRMSSE averaged over every level, by period, on the fast lines (those items), the slow lines (the other items) and the
series with short history (first sale less than 26 weeks or 182 days before the forecast). WM, LGB and Avg as in
Table~\ref{tab:rq-grain}. Shading marks the lowest of the three: \colorbox{wmcell}{world model}, \colorbox{lgbcell}{LightGBM}, \colorbox{avgcell}{average}.}
\label{tab:rq-lines}
\resizebox{\textwidth}{!}{\begin{tabular}{l c ccc ccc ccc}
\toprule
 & Items making & \multicolumn{3}{c}{Fast lines} & \multicolumn{3}{c}{Slow lines} & \multicolumn{3}{c}{Short history}\\
\cmidrule(lr){3-5}\cmidrule(lr){6-8}\cmidrule(lr){9-11}
Data set & 80\% of units & WM & LGB & Avg & WM & LGB & Avg & WM & LGB & Avg\\
\midrule
Analgesics & 17\% & 0.754 & 0.748 & \cellcolor{avgcell}\textbf{0.738} & 0.756 & 0.769 & \cellcolor{avgcell}\textbf{0.748} & \cellcolor{wmcell}\textbf{1.291} & 1.459 & 1.350 \\
Bottled juice & 21\% & 0.594 & 0.606 & \cellcolor{avgcell}\textbf{0.571} & 0.513 & 0.516 & \cellcolor{avgcell}\textbf{0.494} & \cellcolor{wmcell}\textbf{1.303} & 1.393 & 1.312 \\
Cheese & 17\% & 0.505 & 0.507 & \cellcolor{avgcell}\textbf{0.461} & \cellcolor{wmcell}\textbf{0.515} & 0.584 & 0.526 & \cellcolor{wmcell}\textbf{1.222} & 1.411 & 1.284 \\
Frozen entrees & 19\% & 0.327 & 0.310 & \cellcolor{avgcell}\textbf{0.306} & 0.461 & 0.412 & \cellcolor{avgcell}\textbf{0.402} & 1.473 & 1.471 & \cellcolor{avgcell}\textbf{1.336} \\
Soft drinks & 11\% & 0.640 & 0.602 & \cellcolor{avgcell}\textbf{0.583} & 0.604 & 0.651 & \cellcolor{avgcell}\textbf{0.595} & 1.379 & 1.397 & \cellcolor{avgcell}\textbf{1.334} \\
Snacks & 17\% & 0.650 & \cellcolor{lgbcell}\textbf{0.496} & 0.524 & 1.189 & 1.154 & \cellcolor{avgcell}\textbf{1.144} & \cellcolor{wmcell}\textbf{4.204} & 4.565 & 4.328 \\
Frat & 42\% & 0.432 & 0.442 & \cellcolor{avgcell}\textbf{0.414} & 0.659 & 0.626 & \cellcolor{avgcell}\textbf{0.622} & 2.703 & 2.642 & \cellcolor{avgcell}\textbf{2.608} \\
M5 & 37\% & 0.717 & 0.700 & \cellcolor{avgcell}\textbf{0.662} & 0.735 & 0.748 & \cellcolor{avgcell}\textbf{0.702} & \cellcolor{wmcell}\textbf{0.866} & 1.024 & 0.907 \\
Favorita & 42\% & 0.417 & 0.442 & \cellcolor{avgcell}\textbf{0.395} & 0.496 & 0.522 & \cellcolor{avgcell}\textbf{0.472} & 0.539 & 0.667 & \cellcolor{avgcell}\textbf{0.525} \\
\bottomrule
\end{tabular}}
\end{table*}


\subsection{RQ3: seen and unseen data}
\label{sec:rq-unseen}

\textbf{On items it was not trained on, the world model loses almost no accuracy
and LightGBM loses a great deal; on stores it was not trained on, the world model
loses less on three of five data sets and the same on two.} With a fifth of the
items kept out of training, the world model's error on them rises by 0.008 to
0.040, and LightGBM's by 0.058 to 0.313; the difference is clear on all five data
sets tested (Table~\ref{tab:rq-transfer}). LightGBM under-forecasts the kept-out
items, at 66 to 93\% of actual sales; the world model forecasts 89 to 104\%. With a
fifth of the stores kept out, the world model loses less on cheese, snacks and
Frat, and the two lose the same on M5 and Favorita.

\begin{table*}[t]
\centering\footnotesize\setlength{\tabcolsep}{3pt}
\caption{RQ3, seen and unseen data. A fifth of the items, or of the stores, is kept out of training; the cost is the
WRMSSE on that fifth minus the WRMSSE on the same fifth of the same model trained on everything, mean over windows.
The interval is a 90\% bootstrap interval over windows for the world model's cost minus LightGBM's; ``clearly lower''
names the model whose cost is lower when the interval excludes zero. Forecast over actual is on the kept-out series.}
\label{tab:rq-transfer}
\resizebox{\textwidth}{!}{\begin{tabular}{ll cc c l cc}
\toprule
 & & \multicolumn{2}{c}{Cost of not training on them} & & & \multicolumn{2}{c}{Forecast over actual}\\
\cmidrule(lr){3-4}\cmidrule(lr){7-8}
Data set & Kept out & WM & LGB & WM $-$ LGB, 90\% & Clearly lower & WM & LGB\\
\midrule
Cheese & Items & 0.011 & 0.058 & [-0.064, -0.030] & World model & 0.99 & 0.93 \\
Cheese & Stores & 0.001 & 0.047 & [-0.060, -0.032] & World model & 1.05 & 0.94 \\
Snacks & Items & 0.012 & 0.162 & [-0.195, -0.109] & World model & 0.89 & 0.68 \\
Snacks & Stores & 0.001 & 0.103 & [-0.136, -0.068] & World model & 1.00 & 0.72 \\
Frat & Items & 0.008 & 0.313 & [-0.398, -0.248] & World model & 0.95 & 0.66 \\
Frat & Stores & 0.007 & 0.206 & [-0.210, -0.187] & World model & 0.99 & 0.77 \\
M5 & Items & 0.040 & 0.083 & [-0.057, -0.029] & World model & 1.04 & 0.84 \\
M5 & Stores & 0.022 & 0.020 & [-0.001, +0.006] & Neither & 0.99 & 0.93 \\
Favorita & Items & 0.010 & 0.171 & [-0.172, -0.145] & World model & 1.01 & 0.77 \\
Favorita & Stores & 0.020 & 0.017 & [-0.005, +0.015] & Neither & 0.97 & 0.98 \\
\bottomrule
\end{tabular}}
\end{table*}


\subsection{RQ4: transferability}
\label{sec:rq-transfer}

\textbf{The settings carry to another retailer, and training on other retailers'
data does not help either model.} M5's tuned settings carried unchanged to
Favorita give 0.447 overall for the world model, against 0.428 with Favorita's own
settings, and 0.450 for LightGBM, against 0.448 (Table~\ref{tab:rq-retailer});
at item by store the difference is 0.002 for the world model and 0.001 for
LightGBM. Training each model on its data set together with other data sets
lowers the error by no more than 0.007 overall on any data set: by 0.004 for the
world model on analgesics and by 0.007 for LightGBM on frozen entrees. Elsewhere
it leaves the error unchanged or raises it, by as much as 0.141 for the world model on soft drinks and
0.235 for LightGBM on analgesics.

\begin{table*}[t]
\centering\footnotesize\setlength{\tabcolsep}{3pt}
\caption{RQ4, transferability. WRMSSE overall / item by store of each tuned model trained on its own data set and
trained on it together with other data sets (each Dominick's category with the other five and Frat; Frat with the
six Dominick's categories; M5 with Favorita and Favorita with M5). The last row carries M5's tuned settings
unchanged to Favorita, to be read against Favorita's own-data row.}
\label{tab:rq-retailer}
\resizebox{\textwidth}{!}{\begin{tabular}{l cc cc}
\toprule
 & \multicolumn{2}{c}{World model} & \multicolumn{2}{c}{LightGBM}\\
\cmidrule(lr){2-3}\cmidrule(lr){4-5}
Data set & Own data & Pooled & Own data & Pooled\\
\midrule
Analgesics & 0.695 / 0.669 & 0.691 / 0.671 & 0.701 / 0.657 & 0.936 / 0.723 \\
Bottled juice & 0.573 / 0.615 & 0.613 / 0.610 & 0.576 / 0.626 & 0.582 / 0.653 \\
Cheese & 0.462 / 0.639 & 0.462 / 0.632 & 0.512 / 0.618 & 0.573 / 0.650 \\
Frozen entrees & 0.351 / 0.422 & 0.356 / 0.431 & 0.335 / 0.396 & 0.328 / 0.406 \\
Soft drinks & 0.744 / 0.600 & 0.885 / 0.747 & 0.721 / 0.570 & 0.805 / 0.594 \\
Snacks & 1.009 / 0.579 & 1.038 / 0.572 & 0.865 / 0.565 & 0.918 / 0.610 \\
Frat & 0.490 / 0.568 & 0.492 / 0.570 & 0.494 / 0.563 & 0.548 / 0.626 \\
M5 & 0.709 / 0.838 & 0.733 / 0.840 & 0.697 / 0.858 & 0.851 / 0.876 \\
Favorita & 0.428 / 0.690 & 0.448 / 0.699 & 0.448 / 0.722 & 0.459 / 0.738 \\
\midrule
Favorita, M5's settings & 0.447 / 0.692 & & 0.450 / 0.723 & \\
\bottomrule
\end{tabular}}
\end{table*}


\subsection{The broad view}
\label{sec:rq-broad}

Brought together, the four answers give one recommendation and one reason for
it. One forecast made at item by store and summed to every grain, as the
equal-weight average of a tuned world model and tuned gradient-boosted trees,
is the most accurate choice at the grains and on the kinds of line studied
here: it is lowest at store grain and above on eight of nine data sets, at the
item levels on all nine, and on the fast lines and the slow tail on most. The
two models on their own are close on items and stores they were trained on, and
each leads on different data sets. They part company on items and stores with
little or no history, where the world model holds its accuracy and LightGBM
under-forecasts. That is what the world model brings to the average, and where
it fits in the planning cycle: new lines, range changes and new stores, where a
retailer has the least history and the most to lose from a forecast that is too
low. These findings are for the price and promotion levers, the forecast windows
and the data sets studied; Section~\ref{sec:limits} sets out what they do not
cover.


% =====================================================================
\section{Supporting analyses}
\label{sec:support}

Every analysis in this section was rerun on the tuned models of
Table~\ref{tab:final-settings}, on the same rolling windows as
Section~\ref{sec:results}; none of its numbers comes from the settings of the
first version.

\subsection{Every model compared}
\label{sec:sup-comparison}

\textbf{The tuned pair and their average are more accurate than every model not
tuned on the data, on every data set but one.} Table~\ref{tab:sup-comparison}
adds the pretrained models, the Temporal Fusion Transformer and seasonal naive.
TimesFM given the plan is the most accurate of the models not trained on the
data on \supTfmPlanBestUntrained{} data sets. On Frat it is more accurate than
either tuned model alone (\supFratTfmPlan{} against \supFratWm{} for the world
model and \supFratLgb{} for LightGBM) but not than their average
(\supFratAvg{}). Timestral given the plan forecasts far above sales on the
weekly data, where a few extreme upper quantiles dominate the mean of its
output: \supAnaTstPlan{} on analgesics against \supAnaTfmPlan{} for TimesFM given
the plan.

\begin{table*}[t]
\centering\footnotesize\setlength{\tabcolsep}{3pt}
\caption{Every model on the rolling windows: WRMSSE averaged over all levels, mean over windows. WM and LGB: the tuned
world model and LightGBM; Avg: their average; Ens: the three-model ensemble with TimesFM given the plan; TFM+ and
TFM: TimesFM with and without the plan; Chr: Chronos-Bolt; TST+ and TST: Timestral with and without the plan; TFT:
the Temporal Fusion Transformer, with the settings of the first version; SN: seasonal naive. The comparison models
are scored on cleaned prices on Dominick's; on M5, Frat and Favorita they are the runs of the first revision on the
same windows, since nothing in them was tuned. -- marks a model not run on that data set. Shading marks the lowest: \colorbox{wmcell}{world model}, \colorbox{lgbcell}{LightGBM}, \colorbox{avgcell}{average}.}
\label{tab:sup-comparison}
\resizebox{\textwidth}{!}{\begin{tabular}{lccccccccccc}
\toprule
Data set & WM & LGB & Avg & Ens & TFM+ & TFM & Chr & TST+ & TST & TFT & SN\\
\midrule
Analgesics & 0.699 & 0.702 & 0.687 & \textbf{0.676} & 0.723 & 0.804 & 0.846 & 13.892 & 1.341 & 0.817 & 1.374 \\
Bottled juice & 0.575 & 0.581 & 0.551 & \textbf{0.551} & 0.673 & 0.791 & 0.802 & 3.889 & 1.012 & 0.741 & 1.279 \\
Cheese & 0.502 & 0.532 & 0.477 & \textbf{0.471} & 0.612 & 0.727 & 0.765 & 3.437 & 0.770 & 0.600 & 1.187 \\
Frozen entrees & 0.357 & 0.338 & \cellcolor{avgcell}\textbf{0.332} & 0.349 & 0.476 & 0.529 & 0.569 & 1.781 & 0.571 & 0.457 & 0.677 \\
Soft drinks & 0.705 & 0.680 & \cellcolor{avgcell}\textbf{0.650} & 0.655 & 0.770 & 0.909 & 0.976 & 2.281 & 0.858 & 0.815 & 1.001 \\
Snacks & 0.893 & \cellcolor{lgbcell}\textbf{0.778} & 0.790 & 0.795 & 0.855 & 1.079 & 2.316 & 8.666 & 3.277 & 0.959 & 1.399 \\
Frat & 0.490 & 0.494 & 0.469 & \textbf{0.454} & 0.482 & 0.691 & 0.722 & 1.002 & 0.685 & 0.571 & 1.038 \\
M5 & 0.709 & 0.697 & \cellcolor{avgcell}\textbf{0.658} & 0.729 & 0.982 & 0.964 & 1.184 & 1.038 & 0.988 & -- & 1.146 \\
Favorita & 0.428 & 0.448 & \cellcolor{avgcell}\textbf{0.403} & 0.419 & 0.501 & -- & -- & 0.515 & 0.548 & -- & 0.662 \\
\bottomrule
\end{tabular}}
\end{table*}


\subsection{What each member adds to the combination}
\label{sec:sup-ablation}

\textbf{Each tuned model makes the combination more accurate, and TimesFM adds a
little at item by store.} Table~\ref{tab:sup-ablation} drops each member of the
three-model ensemble in turn. Without the world model the ensemble is less
accurate item by store on \supNoWmWorse{} of \supAblationSets{} data sets, and
clearly less accurate in \supPNoWmWorse{} of \supPNoWmN{} window comparisons
(clearly more accurate in \supPNoWmBetter{}); without LightGBM, in
\supPNoLgbWorse{} (clearly more accurate in \supPNoLgbBetter{}). Without TimesFM,
which leaves the average of the two tuned models, it is clearly less accurate item
by store in \supPNoTfmWorse{} comparisons and clearly more accurate in
\supPNoTfmBetter{}. Item by store, then, adding TimesFM with the plan to the
average helps; over all levels the average is in the model confidence set more
often (Section~\ref{sec:rq-grain}). A retailer whose main use is store orders can
add it at the cost of a third model to run. Weights fitted on earlier windows to
the item-by-store error beat equal weights in \supPFitBottomBetter{} of
\supPFitBottomN{} comparisons and lost in \supPFitBottomWorse{}; weights fitted to
the error over all levels beat them in \supPFitAllBetter{} and lost in
\supPFitAllWorse{}.

\begin{table}[t]
\centering\footnotesize\setlength{\tabcolsep}{3pt}
\caption{What each member adds to the three-model ensemble: WRMSSE item by store, mean over the rolling windows, of
the ensemble, of the ensemble without each member (without TimesFM it is the average of the two tuned models), and of
the ensemble with weights fitted on earlier windows. Lowest in bold.}
\label{tab:sup-ablation}
\begin{adjustbox}{max width=\columnwidth}\begin{tabular}{lccccc}
\toprule
Data set & Three & No WM & No LGB & No TFM & Fitted\\
\midrule
Analgesics & 0.655 & 0.655 & 0.664 & 0.655 & \textbf{0.651} \\
Bottled juice & 0.584 & 0.598 & 0.589 & 0.599 & \textbf{0.558} \\
Cheese & 0.599 & 0.609 & 0.617 & 0.605 & \textbf{0.598} \\
Frozen entrees & 0.401 & 0.416 & 0.422 & 0.395 & \textbf{0.336} \\
Soft drinks & \textbf{0.547} & 0.550 & 0.567 & 0.561 & 0.568 \\
Snacks & 0.556 & 0.559 & 0.576 & 0.556 & \textbf{0.531} \\
Frat & \textbf{0.544} & 0.547 & 0.555 & 0.550 & 0.549 \\
M5 & \textbf{0.832} & 0.839 & \textbf{0.832} & 0.836 & 0.849 \\
Favorita & 0.687 & 0.702 & 0.685 & 0.689 & \textbf{0.678} \\
\bottomrule
\end{tabular}\end{adjustbox}
\end{table}


\subsection{Which plan inputs each model uses}
\label{sec:sup-plan}

\begin{table*}[t]
\centering\footnotesize\setlength{\tabcolsep}{4pt}
\caption{Which plan inputs each tuned model uses. Change in WRMSSE item by store, mean over windows, when the price
is set to the regular price (price off) and when every plan input is set to nothing planned (plan off); the change
in forecast units on the rows where the plan had something on; and, for the world model, the largest change when its
history before the last few periods is removed. A positive change means the input was helping.}
\label{tab:sup-perturbation}
\begin{tabular}{l ccc ccc c}
\toprule
 & \multicolumn{3}{c}{World model} & \multicolumn{3}{c}{LightGBM} & World model\\
\cmidrule(lr){2-4}\cmidrule(lr){5-7}\cmidrule(lr){8-8}
Data set & Price off & Plan off & Units, plan rows & Price off & Plan off & Units, plan rows & History cut\\
\midrule
Analgesics & +0.019 & +0.061 & -47\% & +0.037 & +0.075 & -51\% & +0.036 \\
Bottled juice & +0.065 & +0.095 & -53\% & +0.072 & +0.098 & -55\% & +0.102 \\
Cheese & +0.095 & +0.125 & -52\% & +0.136 & +0.161 & -58\% & +0.107 \\
Frozen entrees & -- & -- & -- & -- & -- & -- & -- \\
Soft drinks & -- & -- & -- & -- & -- & -- & -- \\
Snacks & -- & -- & -- & -- & -- & -- & -- \\
Frat & -- & -- & -- & -- & -- & -- & -- \\
M5 & -- & -- & -- & -- & -- & -- & -- \\
Favorita & -- & -- & -- & -- & -- & -- & -- \\
\bottomrule
\end{tabular}
\end{table*}


\subsection{What the world model's carried state does}
\label{sec:sup-mechanism}

\begin{table}[t]
\centering\footnotesize\setlength{\tabcolsep}{3pt}
\caption{The tuned world model's carried state, mean over windows. Timing: change in WRMSSE item by store when the
plan's timing is shuffled within each series. Carry-over: forecast change in periods 1, 2 and 4 after a promotion and a
20\% price cut added in period 1 to series with nothing planned, as a share of the forecast without them. Probe: $R^2$
of log sales after the window from the state before the window and from the state rolled through it (first of the
five copies).}
\label{tab:sup-mechanism}
\resizebox{\columnwidth}{!}{\begin{tabular}{lcccccc}
\toprule
Data set & Timing & Period 1 & Period 2 & Period 4 & Probe before & Probe after\\
\midrule
Analgesics & +16\% & +199\% & +13\% & +11\% & 0.80 & 0.81 \\
Bottled juice & +29\% & +188\% & +6\% & +9\% & 0.86 & 0.88 \\
Cheese & +34\% & +109\% & -4\% & +3\% & 0.91 & 0.92 \\
Frozen entrees & +64\% & +748\% & +54\% & +17\% & 0.80 & 0.84 \\
Soft drinks & +46\% & +327\% & +34\% & +10\% & 0.85 & 0.86 \\
Snacks & +24\% & +129\% & +10\% & +9\% & 0.86 & 0.87 \\
Frat & +52\% & +59\% & -2\% & -1\% & 0.86 & 0.89 \\
M5 & +2\% & -- & -- & -- & 0.64 & 0.65 \\
\bottomrule
\end{tabular}}
\end{table}


\subsection{Decision costs, forecast ranges and reconciliation}
\label{sec:sup-decisions}

\begin{table}[t]
\centering\footnotesize\setlength{\tabcolsep}{3pt}
\caption{Decision costs of the tuned models. Store orders: cost of ordering each forecast's 0.8 quantile, summed over
windows, relative to LightGBM's. Category plan: absolute error in dollars over sales in dollars at category by store
(\%), mean over windows; lowest shaded.}
\label{tab:sup-decisions}
\resizebox{\columnwidth}{!}{\begin{tabular}{lcccccc}
\toprule
 & \multicolumn{3}{c}{Store orders, ratio 0.8} & \multicolumn{3}{c}{Category plan (\%)}\\
\cmidrule(lr){2-4}\cmidrule(lr){5-7}
Data set & WM & LGB & Avg & WM & LGB & Avg\\
\midrule
Analgesics & 1.012 & 1.000 & 0.996 & \cellcolor{wmcell}\textbf{229072263467.5} & 231453772960.2 & 230263018423.3 \\
Bottled juice & 0.983 & 1.000 & 0.965 & \cellcolor{wmcell}\textbf{446696983472.5} & 1188171054110.2 & 817434018801.6 \\
Cheese & 1.077 & 1.000 & 1.006 & \cellcolor{wmcell}\textbf{2041136082896.2} & 2893856365409.1 & 2467496221462.4 \\
Frozen entrees & 1.022 & 1.000 & 0.985 & \cellcolor{wmcell}\textbf{343561014029.6} & 421695443034.5 & 382628228477.2 \\
\bottomrule
\end{tabular}}
\end{table}

\begin{table}[t]
\centering\footnotesize\setlength{\tabcolsep}{3pt}
\caption{Coverage (\%) of each model's 80\% interval, from its own errors in earlier windows, item by store and at
the total, mean over the scored windows.}
\label{tab:sup-intervals}
\resizebox{\columnwidth}{!}{\begin{tabular}{lcccccc}
\toprule
 & \multicolumn{3}{c}{Item by store} & \multicolumn{3}{c}{Total}\\
\cmidrule(lr){2-4}\cmidrule(lr){5-7}
Data set & WM & LGB & Avg & WM & LGB & Avg\\
\midrule
Analgesics & 91 & 91 & 91 & 82 & 82 & 82 \\
Bottled juice & 72 & 73 & 73 & 85 & 98 & 92 \\
Cheese & 74 & 74 & 75 & 75 & 95 & 92 \\
Frozen entrees & 70 & 71 & 71 & 88 & 82 & 90 \\
Soft drinks & 75 & 75 & 76 & 68 & 68 & 68 \\
Snacks & 73 & 73 & 74 & 48 & 75 & 65 \\
Frat & 88 & 87 & 88 & 92 & 62 & 75 \\
M5 & 94 & 94 & 94 & 73 & 80 & 73 \\
Favorita & 81 & 80 & 80 & 65 & 85 & 81 \\
\bottomrule
\end{tabular}}
\end{table}

\begin{table}[t]
\centering\footnotesize\setlength{\tabcolsep}{3pt}
\caption{Reconciliation against summing one forecast: WRMSSE averaged over all levels, mean over windows, of the best
reconciled ETS forecast, of each tuned model combined with ETS above it by weighted least squares, and of the average
summed upward.}
\label{tab:sup-reconcile}
\begin{adjustbox}{max width=\columnwidth}\begin{tabular}{lcccc}
\toprule
Data set & Best ETS & WM + ETS & LGB + ETS & Avg\\
\midrule
Analgesics & 0.836 & 0.903 & 0.933 & \cellcolor{avgcell}\textbf{0.687} \\
Bottled juice & 0.878 & 0.849 & 0.853 & \cellcolor{avgcell}\textbf{0.551} \\
Cheese & 0.742 & 1.430 & 1.215 & \cellcolor{avgcell}\textbf{0.477} \\
Frozen entrees & 0.549 & 0.510 & 0.474 & \cellcolor{avgcell}\textbf{0.332} \\
Soft drinks & \textbf{0.617} & 1.061 & 0.907 & 0.650 \\
Snacks & \textbf{0.698} & 0.979 & 0.805 & 0.790 \\
Frat & 0.693 & 0.499 & \textbf{0.466} & 0.469 \\
M5 & 0.945 & 0.821 & 0.813 & \cellcolor{avgcell}\textbf{0.658} \\
Favorita & 0.570 & 0.434 & 0.462 & \cellcolor{avgcell}\textbf{0.403} \\
\bottomrule
\end{tabular}\end{adjustbox}
\end{table}


% =====================================================================
\section{What was tried and dropped}
\label{sec:dropped}

Most of what was tried did not stay. The searches of Section~\ref{sec:loop} kept
a small share of the changes they ran (Table~\ref{tab:loop}); every change tried
alone in the first round is in Figure~\ref{fig:loopheat}, and the losses tried in
the first revision are in Table~\ref{tab:losses}. Table~\ref{tab:dropped} lists
the larger things the revision tried beyond the searches and did not keep, and
when each becomes worth trying again.

\begin{table*}[t]
\centering
\caption{\textbf{What the revision tried beyond the searches and did not keep.}
Each was written down before it ran.}
\label{tab:dropped}
\footnotesize
\renewcommand{\arraystretch}{1.15}
\begin{tabular}{p{0.25\textwidth}p{0.36\textwidth}p{0.31\textwidth}}
\toprule
\textbf{Tried} & \textbf{What happened} & \textbf{Worth revisiting when} \\
\midrule
Count likelihoods that treat zero weeks on their own (negative binomial, zero-inflated, hurdle), on the world model & Kept on M5 and Favorita (hurdle) and snacks (negative binomial); nothing on the other six data sets & Lines are slower or more often out of stock than here \\
Pretraining the world model's encoder on the latent objective, then fitting the forecast at a lower rate & Kept on snacks only, on the search windows; with nothing kept on eight of nine data sets, Favorita's rerun on its final loss was stopped early & More unlabelled history is available than the forecast needs, as across many categories \\
Training on other retailers' data with each own data set (RQ4) & Helped neither model by more than \rqPoolGainMax{} on any data set, and hurt LightGBM most & The other data come from the same retailer or the same categories \\
The world model alone for series with less than 26 weeks of history & No better than the average on lines with history; worse on snacks & A longer history threshold is set from a retailer's own range data \\
Weights fitted to earlier windows in place of equal weights & Reported in Section~\ref{sec:support} & Many more windows are available to fit them on \\
Timestral in place of TimesFM in the three-model ensemble & Reported in Section~\ref{sec:support} & A retail-trained version is released \\
\bottomrule
\end{tabular}
\end{table*}

The first version of this paper also tried a stock term for the dip after a
promotion, extra loss weight on the periods after a promotion, and a graph model
sharing information between stores; none was kept, and they were not run again on
the tuned models. Their record is in the first version.

% =====================================================================
\section{Discussion}
\label{sec:discussion}

Table~\ref{tab:advice} maps the rule of Section~\ref{sec:rq-broad} to the
decisions a retailer makes from the forecast. Each row is what the tuned models
did on the nine data sets studied, within the scope of
Section~\ref{sec:limits}.

\begin{table*}[t]
\centering
\caption{The decisions a retailer makes from one item-by-store forecast, and what
the tuned models did at each on the nine data sets studied.}
\label{tab:advice}
\footnotesize
\renewcommand{\arraystretch}{1.2}
\begin{tabular}{p{0.21\textwidth}p{0.17\textwidth}p{0.18\textwidth}p{0.36\textwidth}}
\toprule
\textbf{Decision} & \textbf{Grain} & \textbf{Forecast to use} & \textbf{What the tuned models did} \\
\midrule
Store orders and inventory & Item by store, by day or week & The average & Lowest error at the item levels on all \rqAvgBestItemWord{} data sets; clearly less accurate than either model in at most \rqAvgItemWorseMax{} of \rqPAvgWmItemn{} paired comparisons \\
Promotion and supplier funding & Item across stores, by week and over the promotion block & The average & Lowest at the item levels on all \rqAvgBestItemWord{}; over the whole block lowest on \rqAvgBestBlockWord{} and second on the other four; LightGBM alone drifts in total on M5, Favorita and soft drinks \\
Category performance & Category by store and category, by week & The average & Lowest at store grain and above on \rqAvgBestCoarseWord{} of \rqNineWord{}; LightGBM lowest on snacks \\
Financial plan & Chain total & The average & As for category performance; the total is the sum of the same item-by-store forecast, so it agrees with every plan below it \\
New items and range changes & Item by store, items not yet trained on & The world model alone & Clearly more accurate than the average on all \rqKeptSets{} data sets tested; LightGBM forecast \rqFaLgbItemLo{} to \rqFaLgbItemHi{}\% of actual sales on these items \\
New stores & Item by store, stores not yet trained on & The average & The average and the world model each clearly ahead on \rqStoreAvgBeatsWm{} data set and level on \rqStoreLevel{}; LightGBM clearly behind the world model on \rqStoreWmBeatsLgb{} of \rqKeptSets{} \\
Choosing for another retailer & Any & Rerun the comparison & The order of models is a finding on these data; every comparison, including those that lost here, should be rerun on the retailer's own data \\
\bottomrule
\end{tabular}
\end{table*}

\paragraph{Why the average works}
The two tuned models are close in accuracy and wrong in different places. Across
the nine data sets, each is clearly more accurate than the other in a similar
number of paired comparisons (\rqPWmLgbBetter{} for the world model, \rqPWmLgbWorse{} for LightGBM), and the
data sets on which each leads differ. An equal-weight average of two forecasts
whose errors differ is more accurate than either on most series, and no weights
have to be estimated, which the forecast-combination literature finds is hard to
beat \cite{smithwallis2009,wang2023}.

\paragraph{What the world model brings}
The world model's contribution is clearest where the forecast has least to go
on. On items it was not trained on it lost at most \rqCostWmItemHi{} in scaled error, against
up to \rqCostLgbItemHi{} for LightGBM, and forecast within \rqWmFaItemOff\% of actual sales; LightGBM
forecast as little as \rqLgbFaItemLoFrac{}. It reads an item's recent sales into a state
and carries that state forward under the plan, so it does not need the item in
its training data, whereas LightGBM takes the item as one of its labels and has
learned nothing for an item it was not trained on. The same carried state keeps its totals over the
promotion block closer to actual sales than LightGBM's on M5, Favorita and soft
drinks.

\paragraph{What the trees bring}
LightGBM is the more accurate of the two on frozen entrees and snacks at every
grain, and its block totals are closer on analgesics, bottled juice and snacks.
It needs no pretraining and trains in minutes, and on lines with history it is
as accurate as the world model on most data sets. On snacks, where the world
model alone is weakest (\rqSnaWm{} overall against \rqSnaLgb{} for LightGBM), the average
stays close to LightGBM at \rqSnaAvg{}.

% =====================================================================
\section{Limitations and threats to validity}
\label{sec:limits}
\label{sec:limitations}

\paragraph{Scope: two of the four Ps, on public data}
The paper addresses price and promotion. It forecasts items newly added to a
range but does not decide product (which items to range) or place (space,
distribution), and it does not solve the four Ps together; a unified forecast across them does not exist in
this work, and one could change which model works best at each grain. The data
are four public data sets with their own plan inputs and constraints. A single
retailer's data, with its own levers, stock and space records, could give a
different ordering of models, and the grain comparison should be repeated on
it before a choice is made.

\paragraph{Adapting the findings}
The rule of Section~\ref{sec:rq-broad} and the ordering of models behind it hold
on the nine data sets studied. They are not a ranking of methods for retail in
general. A retailer or researcher adopting the framework should rerun every
comparison on its own data, including the models and combinations that lost
here, since data with other levers, lines and lengths of history can change the
order. The code and settings for every comparison are released for that purpose.

\paragraph{Forecast horizon}
Every forecast in the study runs 8 weeks ahead on the weekly data and 28 days
ahead on the daily data. Whether the models' accuracy, and the order between
them, holds over the 26- and 52-week horizons of range, space and supplier
reviews was not tested: a longer forecast may degrade, and may degrade
differently for each model. Testing it needs a retailer's own plans and range
decisions to score against and is left for separate work.

\paragraph{New items}
The kept-out items of the transfer test were not trained on, but each model read
their recent sales before forecasting them. An item with no sales at all, on its
first day in range, was not tested; the new-item part of the rule covers items
the models have not been trained on, not items with no history. That part of
the rule was drawn from the results, after a wider rule written down in advance
held only in part (Section~\ref{sec:rq-broad}).

\paragraph{Windows and seeds}
Four to six windows per data set is a thin basis for a bootstrap over windows
and for a model confidence set: the intervals are approximate, the set has
little power, and a model kept in it has not been shown to be as good as the
best. Adjacent windows share training history and season, so they carry less
information than their number suggests; the item bootstrap within a window
treats the items as the sample and does not cover variation between windows.
Averaging five fits with different seeds removes much of the noise between runs
that the searches measured (Section~\ref{sec:loopnoise}), but both models on M5
are single fits.

\paragraph{Analyses added after the first results}
The analyses by line type, the whole-block view, the forecast by history and the
supporting reruns were added after the first scores of the tuned models were
read. Each was written down before it ran, but which to run was chosen knowing
the first results, and the new-item part of the recommended rule was drawn from
the results (Section~\ref{sec:rq-broad}).

\paragraph{Comparison models}
Only the world model and LightGBM were tuned; the Temporal Fusion Transformer
runs with the settings of the first version, and the pretrained models as
published. The comparison between the two tuned models is therefore the fair one,
and the others are reference points. TFT was not run on M5 or Favorita, and
Chronos-Bolt and TimesFM from history alone not on Favorita. TimesFM with the
plan fits its linear plan model within groups of 2,048 series. The pretrained
models' point forecast was fixed in advance as the mean of their output, which
for Timestral let a few extreme upper quantiles dominate on weekly data.

\paragraph{What lift means}
Merchants promote in weeks they expect to sell well, so part of a measured lift
is good timing. The lift is what the forecast attributes to the plan, not what
the plan caused. Our companion audit~\cite{paper2} measures the gap in generated store
data where the true lifts are known. Plan responses in this paper, including the
perturbation tests, are forecasts, not causal estimates.

\paragraph{What the data do not record}
None of the data sets records stock, so a sell-out looks like low demand; or
shelf space. Favorita records no prices. The data were collected between 1989
and 2017 and their plan inputs differ, so the results compare models within each
data set. The decision costs are simulations under stated costs.

% =====================================================================
\section{Conclusion}
\label{sec:conclusion}

A retailer makes decisions from the same demand at different grains, and on
these data no single model was the most accurate at every grain. One forecast
made item by store and summed upward, as the equal-weight average of a tuned
world model and tuned gradient-boosted trees, was the most accurate choice at
store grain and above on \rqAvgBestCoarseWord{} of \rqNineWord{} data sets and at the item levels on all
\rqAvgBestItemWord{}, on the fast lines and on the slow tail. For items the models had not been
trained on, the world model alone was more accurate than the average on every
data set tested, because LightGBM under-forecast them. Settings carried from one
retailer's data to another's, and training on other retailers' data did not
help. The rule that follows, the average everywhere and the world model alone
for new items, is a finding on these data for forecasts of 8 weeks and 28 days
under price and promotion plans. The comparison behind it, run over several
windows with paired tests and every losing model reported, is what we propose a
retailer reruns on its own data, and over its own planning horizons, before
adopting it.

\section{Reproducibility and compute}
\label{sec:repro}

\paragraph{Records}
All code, configuration files and tests are in the repository. A run refuses to
start on uncommitted code and writes a manifest holding the code commit, the
full configuration and its checksum, the seed, checksums of every data file
read, package versions, the machine and run times. The scores and manifests of
the final, transfer and pooled runs, every try of the searches, and the
analyses of the tuned models are committed under \nolinkurl{results/final/},
\nolinkurl{results/ftransfer/}, \nolinkurl{results/pooled/},
\nolinkurl{results/research/} and \nolinkurl{results/analysis_tuned/}; the plan
and every dated entry are in \nolinkurl{docs/PAPER1_REVISION_PLAN.md}.
\texttt{make} in \texttt{paper/} rebuilds every figure, table and quoted number
from these files (\buildSeconds{} seconds from a fresh clone of the repository,
\buildDate). Forecast files and model checkpoints are kept outside the
repository on the results volume, with their checksums recorded.

\paragraph{Compute}
The first version ran on one Google Colab machine with an NVIDIA \computeGpu{}
(Table~\ref{tab:compute}: \computeRuns{} timed runs, \computeHours{}
machine-hours). The revision ran on rented machines: the searches, final runs,
transfer tests and pooled training on GPUs from Vast, managed and analysed on
Modal. Table~\ref{tab:compute-revision} gives the run time recorded in each of
the \computeRevRuns{} run folders, \computeRevHours{} hours in all, most on RTX
4090 GPUs; because a resumed run records only its last stretch, this is a floor.
The machine time paid for on Vast was \computeVastGpuHours{} GPU-hours
(US\$\computeVastGpuUsd{}), and the revision cost US\$\computeTotalUsd{} in all
from 26 October to 12 November 2026: US\$\computeVastUsd{} on Vast, including
storage and data transfer, and US\$\computeModalUsd{} on Modal. The total
includes every run made, the searches' screens, repeats after machine failures
and runs lost when a machine stopped, not only the runs reported.

\begin{table}[t]
\centering
\caption{\textbf{Run times of the final and comparison runs, from their manifests.} Wall-clock minutes per
run (data loading, training, forecasting and scoring) and hours in total, on one Google Colab machine with an NVIDIA Tesla T4 (LightGBM uses its CPU). Dominick's:
six categories together; a Dominick's world-model run averages three fits. Two M5 runs resumed after the
Colab runtime dropped record only the resumed part (the shortest values in their rows). Not counted: runs saved before
run times were recorded, the development study, and TimesFM on Dominick's and Frat.}
\label{tab:compute}
\footnotesize
\setlength{\tabcolsep}{3.5pt}
\begin{tabular}{llrrlr}
\toprule
\textbf{Data} & \textbf{Model} & \textbf{Runs} & \textbf{Median} & \textbf{Range} & \textbf{Hours} \\
\midrule
M5 & World model & 9 & 101.5 & 37.3 to 103.2 & 14.0 \\
M5 & Lift model & 9 & 88.7 & 10.3 to 94.6 & 12.2 \\
M5 & LightGBM & 9 & 78.1 & 76.5 to 79.9 & 11.7 \\
M5 & Graph model & 12 & 7.2 & 7.2 to 7.5 & 1.5 \\
M5 & TimesFM & 8 & 3.9 & 3.8 to 4.1 & 0.5 \\
M5 & Seasonal naive & 2 & 1.3 & 1.2 to 1.3 & 0.0 \\
Dominick's & World model & 18 & 9.2 & 9.0 to 10.6 & 2.8 \\
Dominick's & LightGBM & 18 & 4.6 & 3.1 to 11.5 & 1.7 \\
Frat & World model & 3 & 2.2 & 2.2 to 2.2 & 0.1 \\
Frat & Lift model & 3 & 2.0 & 1.9 to 2.0 & 0.1 \\
Frat & LightGBM & 3 & 1.5 & 1.5 to 1.6 & 0.1 \\
\midrule
\multicolumn{5}{l}{All 94 timed runs} & 44.7 \\
\bottomrule
\end{tabular}
\end{table}

\begin{table}[t]
\centering
\caption{\textbf{Run time of the \revn{}, from each run's manifest.} Runs counts every run folder, the pieces of
split runs included; hours is their summed wall-clock time; machines are the GPUs or CPU hosts the runs used, with the
number of runs on each. A run resumed from a partway save records only its last stretch, so the hours are a floor;
the machine time paid for is in the text.}
\label{tab:compute-revision}
\footnotesize
\setlength{\tabcolsep}{3pt}
\begin{tabular}{p{0.36\columnwidth}rrp{0.32\columnwidth}}
\toprule
\textbf{Study} & \textbf{Runs} & \textbf{Hours} & \textbf{Machines} \\
\midrule
Searches of both models' settings & 3,958 & 4,305 & RTX 4090 (1182), RTX 5090 (54), L4 (868), CPU (1288), other GPUs (566) \\
Final runs on the rolling windows & 245 & 41 & RTX 4090 (220), RTX 5090 (16), other GPUs (9) \\
Transfer and new-retailer tests & 358 & 150 & RTX 4090 (252), RTX 5090 (85), other GPUs (21) \\
Training on other retailers' data & 249 & 176 & RTX 4090 (216), RTX 5090 (17), other GPUs (16) \\
Comparison models on cleaned prices & 60 & 8 & RTX 4090 (48), other GPUs (12) \\
\midrule
All & 4,870 & 4,679 & \\
\bottomrule
\end{tabular}
\end{table}


\section{Data, code and licence}
\label{sec:licence}

\paragraph{Data}
M5 is available from the competition's Kaggle page~\cite{m5accuracy};
Dominick's from the Kilts Center for Marketing, University of Chicago Booth
School of Business~\cite{dominicks}; Breakfast at the Frat from
dunnhumby~\cite{frat}; Corporaci\'on Favorita from its Kaggle
competition~\cite{favorita}. Each is obtained from its provider under the
provider's own terms. No data is redistributed with the code.

\paragraph{Code}
The code is at
\url{https://github.com/karthik-iyer-au/retail-world-models}; the version used
for this paper is archived as release v1.0.0 on Zenodo,
\href{https://doi.org/10.5281/zenodo.23448056}{doi:10.5281/zenodo.23448056}.

\paragraph{Copyright and licence}
Copyright \copyright{} 2026 the authors. The code is licensed under the
PolyForm Noncommercial License~1.0.0 and the manuscript text, figures and tables
under Creative Commons Attribution-NonCommercial~4.0 International
(CC~BY-NC~4.0). Industry and commercial use of the code, the manuscript or any
part of them, including use inside a company, is prohibited unless approved in
writing by the authors.

\section{Author contributions}
\label{sec:contributions}
Karthik Iyer conceived and scoped the work, designed the study, wrote the code,
ran the experiments and analyses, and wrote the paper, with AI assistance as
declared in \S\ref{sec:ai}. Madan Dabbeeru widened the range of methods
considered for each research question, and reviewed the work and the
manuscript.

\section{Practice background and independence}
\label{sec:background}
The first author has worked in retail for fifteen years, across continents, in
grocery, department stores, fashion, electronics and online marketplaces. No
data, method, code or other intellectual property from those roles or
engagements is used in this work, and no employer or client data are used: the
study is built from the ground up on public data sets. What the authors carry
from practice is the set of questions retailers ask of a forecast, which shaped
the research questions in Section~\ref{sec:framework}. The first author's
research focus is world models; the second author widened the scientific
options considered for answering each question.

\section{Declaration of generative AI and AI-assisted technologies}
\label{sec:ai}
The research questions, scope, study design, choice of data and methods, and
interpretation of results are the authors' own. Claude (Anthropic) was used as an
assistant to write and test code, run analyses, produce figures and tables, and
draft and edit text. The authors directed this work throughout, reviewed every
result, and corrected and revised the code, analyses and text over many
iterations, and take full responsibility for the content of the publication.

\begin{thebibliography}{99}
\footnotesize
\setlength{\itemsep}{1.5pt}\setlength{\parskip}{0pt}

\bibitem{m5accuracy} S.~Makridakis, E.~Spiliotis, and V.~Assimakopoulos. M5 accuracy competition: Results, findings, and conclusions. \emph{International Journal of Forecasting}, 38(4):1346--1364, 2022.

\bibitem{m5uncertainty} S.~Makridakis, E.~Spiliotis, V.~Assimakopoulos, Z.~Chen, A.~Gaba, I.~Tsetlin, and R.~L. Winkler. The M5 uncertainty competition: Results, findings and conclusions. \emph{International Journal of Forecasting}, 38(4):1365--1385, 2022.

\bibitem{theodorou2022} E.~Theodorou, S.~Wang, Y.~Kang, E.~Spiliotis, S.~Makridakis, and V.~Assimakopoulos. Exploring the representativeness of the M5 competition data. \emph{International Journal of Forecasting}, 38(4):1500--1506, 2022.

\bibitem{inandjung} Y.~In and J.-Y. Jung. Simple averaging of direct and recursive forecasts via partial pooling using machine learning. \emph{International Journal of Forecasting}, 38(4):1386--1399, 2022.

\bibitem{fildes2022} R.~Fildes, S.~Ma, and S.~Kolassa. Retail forecasting: Research and practice. \emph{International Journal of Forecasting}, 38(4):1283--1318, 2022.

\bibitem{croston} J.~D. Croston. Forecasting and stock control for intermittent demands. \emph{Operational Research Quarterly}, 23(3):289--303, 1972.

\bibitem{syntetos2005} A.~A. Syntetos and J.~E. Boylan. The accuracy of intermittent demand estimates. \emph{International Journal of Forecasting}, 21(2):303--314, 2005.

\bibitem{kolassa2016} S.~Kolassa. Evaluating predictive count data distributions in retail sales forecasting. \emph{International Journal of Forecasting}, 32(3):788--803, 2016.

\bibitem{jorgensen1987} B.~J{\o}rgensen. Exponential dispersion models. \emph{Journal of the Royal Statistical Society, Series B}, 49(2):127--145, 1987.

\bibitem{promocast} L.~G. Cooper, P.~Baron, W.~Levy, M.~Swisher, and P.~Gogos. PromoCast: A new forecasting method for promotion planning. \emph{Marketing Science}, 18(3):301--316, 1999.

\bibitem{vanheerde2000} H.~J. van~Heerde, P.~S.~H. Leeflang, and D.~R. Wittink. The estimation of pre- and postpromotion dips with store-level scanner data. \emph{Journal of Marketing Research}, 37(3):383--395, 2000.

\bibitem{mace2004} S.~Mac\'e and S.~A. Neslin. The determinants of pre- and postpromotion dips in sales of frequently purchased goods. \emph{Journal of Marketing Research}, 41(3):339--350, 2004.

\bibitem{ma2016} S.~Ma, R.~Fildes, and T.~Huang. Demand forecasting with high dimensional data: The case of SKU retail sales forecasting with intra- and inter-category promotional information. \emph{European Journal of Operational Research}, 249(1):245--257, 2016.
\bibitem{vandonselaar2016} K.~H. van Donselaar, J.~Peters, A.~de Jong, and R.~A.~C.~M. Broekmeulen. Analysis and forecasting of demand during promotions for perishable items. \emph{International Journal of Production Economics}, 172:65--75, 2016.
\bibitem{babai2014} M.~Z. Babai, A.~Syntetos, and R.~Teunter. Intermittent demand forecasting: An empirical study on accuracy and the risk of obsolescence. \emph{International Journal of Production Economics}, 157:212--219, 2014.
\bibitem{babai2019} M.~Z. Babai, Y.~Dallery, S.~Boubaker, and R.~Kalai. A new method to forecast intermittent demand in the presence of inventory obsolescence. \emph{International Journal of Production Economics}, 209:30--41, 2019.
\bibitem{kourentzes2020} N.~Kourentzes, J.~R. Trapero, and D.~K. Barrow. Optimising forecasting models for inventory planning. \emph{International Journal of Production Economics}, 225:107597, 2020.

\bibitem{pepsico2026} A.~Llenas, E.~Salazar-Trevi\~{n}o, F.~Leskovar, F.~Pons-Llopis, F.~Todeschini, D.~Gowda, D.~Bofill, L.~Anish, and M.~Cleavinger. PepsiCo deploys AI-driven pricing and promotion optimization at scale. \emph{INFORMS Journal on Applied Analytics}, 2026. doi:10.1287/inte.2025.0302.

\bibitem{ma2025} B.~J. Ma, I.~Jackson, M.~Huang, S.~Villegas, and J.~Macias-Aguayo. A data-driven and context-aware approach for demand forecasting in the beverage industry. \emph{International Journal of Logistics Research and Applications}, 2025.

\bibitem{gupta1988} S.~Gupta. Impact of sales promotions on when, what, and how much to buy. \emph{Journal of Marketing Research}, 25(4):342--355, 1988.

\bibitem{vanheerde2003} H.~J. van~Heerde, S.~Gupta, and D.~R. Wittink. Is 75\% of the sales promotion bump due to brand switching? No, only 33\% is. \emph{Journal of Marketing Research}, 40(4):481--491, 2003.

\bibitem{bell1975} D.~E. Bell, R.~L. Keeney, and J.~D.~C. Little. A market share theorem. \emph{Journal of Marketing Research}, 12(2):136--141, 1975.

\bibitem{cooper1988} L.~G. Cooper and M.~Nakanishi. \emph{Market-Share Analysis: Evaluating Competitive Marketing Effectiveness}. Kluwer Academic Publishers, Boston, 1988.

\bibitem{hoch1995} S.~J. Hoch, B.-D. Kim, A.~L. Montgomery, and P.~E. Rossi. Determinants of store-level price elasticity. \emph{Journal of Marketing Research}, 32(1):17--29, 1995.

\bibitem{dominicks} James M. Kilts Center for Marketing. Dominick's dataset. University of Chicago Booth School of Business. \url{https://www.chicagobooth.edu/research/kilts/research-data/dominicks}.

\bibitem{frat} dunnhumby. Breakfast at the Frat. dunnhumby source files. \url{https://www.dunnhumby.com/source-files/}.

\bibitem{athanasopoulos2009} G.~Athanasopoulos, R.~A. Ahmed, and R.~J. Hyndman. Hierarchical forecasts for Australian domestic tourism. \emph{International Journal of Forecasting}, 25(1):146--166, 2009.

\bibitem{athanasopoulos2017} G.~Athanasopoulos, R.~J. Hyndman, N.~Kourentzes, and F.~Petropoulos. Forecasting with temporal hierarchies. \emph{European Journal of Operational Research}, 262(1):60--74, 2017.

\bibitem{kourentzes2019crosstemporal} N.~Kourentzes and G.~Athanasopoulos. Cross-temporal coherent forecasts for Australian tourism. \emph{Annals of Tourism Research}, 75:393--409, 2019.

\bibitem{hu2019newproduct} K.~Hu, J.~Acimovic, F.~Erize, D.~J. Thomas, and J.~A. Van~Mieghem. Forecasting new product life cycle curves: practical approach and empirical analysis. \emph{Manufacturing \& Service Operations Management}, 21(1):66--85, 2019.

\bibitem{vansteenbergen2020} R.~M. van~Steenbergen and M.~R.~K. Mes. Forecasting demand profiles of new products. \emph{Decision Support Systems}, 139:113401, 2020.

\bibitem{hyndman2011} R.~J. Hyndman, R.~A. Ahmed, G.~Athanasopoulos, and H.~L. Shang. Optimal combination forecasts for hierarchical time series. \emph{Computational Statistics \& Data Analysis}, 55(9):2579--2589, 2011.

\bibitem{mint} S.~L. Wickramasuriya, G.~Athanasopoulos, and R.~J. Hyndman. Optimal forecast reconciliation for hierarchical and grouped time series through trace minimization. \emph{Journal of the American Statistical Association}, 114(526):804--819, 2019.

\bibitem{hyndmankoehler} R.~J. Hyndman and A.~B. Koehler. Another look at measures of forecast accuracy. \emph{International Journal of Forecasting}, 22(4):679--688, 2006.

\bibitem{tashman2000} L.~J. Tashman. Out-of-sample tests of forecasting accuracy: An analysis and review. \emph{International Journal of Forecasting}, 16(4):437--450, 2000.

\bibitem{montero2021} P.~Montero-Manso and R.~J. Hyndman. Principles and algorithms for forecasting groups of time series: Locality and globality. \emph{International Journal of Forecasting}, 37(4):1632--1653, 2021.

\bibitem{lightgbm} G.~Ke, Q.~Meng, T.~Finley, T.~Wang, W.~Chen, W.~Ma, Q.~Ye, and T.-Y. Liu. LightGBM: A highly efficient gradient boosting decision tree. \emph{Advances in Neural Information Processing Systems 30}, 2017.

\bibitem{januschowski2022} T.~Januschowski, Y.~Wang, K.~Torkkola, T.~Erkkil\"a, H.~Hasson, and J.~Gasthaus. Forecasting with trees. \emph{International Journal of Forecasting}, 38(4):1473--1481, 2022.

\bibitem{nbeats} B.~N. Oreshkin, D.~Carpov, N.~Chapados, and Y.~Bengio. N-BEATS: Neural basis expansion analysis for interpretable time series forecasting. \emph{International Conference on Learning Representations (ICLR)}, 2020.

\bibitem{deepar} D.~Salinas, V.~Flunkert, J.~Gasthaus, and T.~Januschowski. DeepAR: Probabilistic forecasting with autoregressive recurrent networks. \emph{International Journal of Forecasting}, 36(3):1181--1191, 2020.

\bibitem{tft} B.~Lim, S.~\"O. Ar{\i}k, N.~Loeff, and T.~Pfister. Temporal fusion transformers for interpretable multi-horizon time series forecasting. \emph{International Journal of Forecasting}, 37(4):1748--1764, 2021.

\bibitem{batesgranger} J.~M. Bates and C.~W.~J. Granger. The combination of forecasts. \emph{Operational Research Quarterly}, 20(4):451--468, 1969.

\bibitem{clemen1989} R.~T. Clemen. Combining forecasts: A review and annotated bibliography. \emph{International Journal of Forecasting}, 5(4):559--583, 1989.

\bibitem{graphsage} W.~L. Hamilton, R.~Ying, and J.~Leskovec. Inductive representation learning on large graphs. \emph{Advances in Neural Information Processing Systems 30}, 2017.

\bibitem{veit2017} A.~Veit, S.~Belongie, and T.~Karaletsos. Conditional similarity networks. \emph{Proceedings of the IEEE Conference on Computer Vision and Pattern Recognition (CVPR)}, pages 830--838, 2017.

\bibitem{timesfm} A.~Das, W.~Kong, R.~Sen, and Y.~Zhou. A decoder-only foundation model for time-series forecasting. \emph{Proceedings of the 41st International Conference on Machine Learning (ICML)}, 2024.

\bibitem{timesfm25} Google Research. TimesFM repository: TimesFM 2.5 (200M) and covariate support through XReg. \url{https://github.com/google-research/timesfm}, 2025; checkpoint \texttt{google/timesfm-2.5-200m-pytorch}, package \texttt{timesfm} 2.0.2.

\bibitem{chronos} A.~F. Ansari, L.~Stella, C.~Turkmen, X.~Zhang, P.~Mercado, H.~Shen, O.~Shchur, S.~S. Rangapuram, S.~Pineda~Arango, S.~Kapoor, J.~Zschiegner, D.~C. Maddix, H.~Wang, M.~W. Mahoney, K.~Torkkola, A.~G. Wilson, M.~Bohlke-Schneider, and Y.~Wang. Chronos: Learning the language of time series. \emph{Transactions on Machine Learning Research}, 2024.

\bibitem{lagllama} K.~Rasul, A.~Ashok, A.~R. Williams, H.~Ghonia, R.~Bhagwatkar, A.~Khorasani, M.~J. Darvishi~Bayazi, G.~Adamopoulos, R.~Riachi, N.~Hassen, M.~Bilo\v{s}, S.~Garg, A.~Schneider, N.~Chapados, A.~Drouin, V.~Zantedeschi, Y.~Nevmyvaka, and I.~Rish. Lag-Llama: Towards foundation models for probabilistic time series forecasting. arXiv:2310.08278, 2023.

\bibitem{moirai} G.~Woo, C.~Liu, A.~Kumar, C.~Xiong, S.~Savarese, and D.~Sahoo. Unified training of universal time series forecasting transformers. \emph{Proceedings of the 41st International Conference on Machine Learning}, PMLR 235:53140--53164, 2024.

\bibitem{timegpt} A.~Garza, C.~Challu, and M.~Mergenthaler-Canseco. TimeGPT-1. arXiv:2310.03589, 2023.

\bibitem{worldmodels} D.~Ha and J.~Schmidhuber. Recurrent world models facilitate policy evolution. \emph{Advances in Neural Information Processing Systems 31}, 2018.

\bibitem{planet} D.~Hafner, T.~Lillicrap, I.~Fischer, R.~Villegas, D.~Ha, H.~Lee, and J.~Davidson. Learning latent dynamics for planning from pixels. \emph{Proceedings of the 36th International Conference on Machine Learning}, PMLR 97:2555--2565, 2019.

\bibitem{dreamerv3} D.~Hafner, J.~Pasukonis, J.~Ba, and T.~Lillicrap. Mastering diverse domains through world models. arXiv:2301.04104, 2023.

\bibitem{lecun2022} Y.~LeCun. A path towards autonomous machine intelligence. OpenReview preprint, 2022.

\bibitem{ijepa} M.~Assran, Q.~Duval, I.~Misra, P.~Bojanowski, P.~Vincent, M.~Rabbat, Y.~LeCun, and N.~Ballas. Self-supervised learning from images with a joint-embedding predictive architecture. \emph{Proceedings of the IEEE/CVF Conference on Computer Vision and Pattern Recognition (CVPR)}, pages 15619--15629, 2023.

\bibitem{vjepa2} M.~Assran, A.~Bardes, D.~Fan, Q.~Garrido, R.~Howes, M.~Komeili, M.~Muckley, A.~Rizvi, C.~Roberts, K.~Sinha, A.~Zholus, S.~Arnaud, A.~Gejji, A.~Martin, F.~R. Hogan, D.~Dugas, P.~Bojanowski, V.~Khalidov, P.~Labatut, F.~Massa, M.~Szafraniec, K.~Krishnakumar, Y.~Li, X.~Ma, S.~Chandar, F.~Meier, Y.~LeCun, M.~Rabbat, and N.~Ballas. V-JEPA 2: Self-supervised video models enable understanding, prediction and planning. arXiv:2506.09985, 2025.

\bibitem{byol} J.-B. Grill, F.~Strub, F.~Altch\'e, C.~Tallec, P.~H. Richemond, E.~Buchatskaya, C.~Doersch, B.~Avila~Pires, Z.~D. Guo, M.~Gheshlaghi~Azar, B.~Piot, K.~Kavukcuoglu, R.~Munos, and M.~Valko. Bootstrap your own latent: A new approach to self-supervised learning. \emph{Advances in Neural Information Processing Systems 33}, pages 21271--21284, 2020.

\bibitem{vicreg} A.~Bardes, J.~Ponce, and Y.~LeCun. VICReg: Variance-invariance-covariance regularization for self-supervised learning. \emph{International Conference on Learning Representations (ICLR)}, 2022.

\bibitem{lejepa} R.~Balestriero and Y.~LeCun. LeJEPA: Provable and scalable self-supervised learning without the heuristics. arXiv:2511.08544, 2025.

\bibitem{leworldmodel} L.~Maes, Q.~Le~Lidec, D.~Scieur, Y.~LeCun, and R.~Balestriero. LeWorldModel: Stable end-to-end joint-embedding predictive architecture from pixels. arXiv:2603.19312, 2026.

\bibitem{tdmpc2} N.~Hansen, H.~Su, and X.~Wang. TD-MPC2: Scalable, robust world models for continuous control. \emph{International Conference on Learning Representations (ICLR)}, 2024.

\bibitem{dwm} Y.-G. Zhang, T.~Du, Q.~Zhang, and Y.~Wang. DWM: Separating world effects from actions in latent world models. arXiv:2607.18715, 2026.

\bibitem{adwm} J.~Qiu, Z.~Chen, H.~Cao, J.~Shi, J.~Huo, and Y.~Gao. AD-WM: Action-discriminative world models for counterfactual model predictive control. arXiv:2609.30264, 2026.

\bibitem{stablewm} L.~Maes, Q.~Le~Lidec, L.~Facury, N.~Massaudi, A.~Chaurasia, F.~Capuano, R.~Gao, T.~Gillin, D.~Haramati, D.~Scieur, Y.~LeCun, and R.~Balestriero. stable-worldmodel: A platform for reproducible world modeling research and evaluation. arXiv:2605.21800, 2026.

\bibitem{worldts} Y.~Zhu, X.~Qiu, H.~Cheng, W.~Shen, C.~Guo, B.~Yang, J.~Hu, and C.~S. Jensen. WorldTS: World modeling for multimodal covariate-aware time series forecasting. arXiv:2609.31162, 2026.

\bibitem{mjepa} \ifjournal Anonymised for review. M-JEPA: joint-embedding predictive pretraining for molecular property prediction. Code and evaluation protocol, 2026.\else K.~Iyer. M-JEPA: joint-embedding predictive pretraining for molecular property prediction, code and evaluation protocol. \url{https://github.com/Kar488/M-JEPA}.\fi

\bibitem{gru} K.~Cho, B.~van~Merri\"enboer, C.~Gulcehre, D.~Bahdanau, F.~Bougares, H.~Schwenk, and Y.~Bengio. Learning phrase representations using RNN encoder--decoder for statistical machine translation. \emph{Proceedings of EMNLP}, 2014.

\bibitem{transformer} A.~Vaswani, N.~Shazeer, N.~Parmar, J.~Uszkoreit, L.~Jones, A.~N. Gomez, {\L}.~Kaiser, and I.~Polosukhin. Attention is all you need. \emph{Advances in Neural Information Processing Systems 30}, 2017.

\bibitem{layernorm} J.~L. Ba, J.~R. Kiros, and G.~E. Hinton. Layer normalization. arXiv:1607.06450, 2016.

\bibitem{dropout} N.~Srivastava, G.~Hinton, A.~Krizhevsky, I.~Sutskever, and R.~Salakhutdinov. Dropout: A simple way to prevent neural networks from overfitting. \emph{Journal of Machine Learning Research}, 15(56):1929--1958, 2014.

\bibitem{adamw} I.~Loshchilov and F.~Hutter. Decoupled weight decay regularization. \emph{International Conference on Learning Representations (ICLR)}, 2019.

\bibitem{paper2} \ifjournal Anonymised for review. Auditing world-model explanations in retail forecasting: faithfulness, accuracy and causal validity. Manuscript, 2026.\else K.~Iyer. Auditing world-model explanations in retail forecasting: faithfulness, accuracy and causal validity. Companion paper, \url{https://github.com/karthik-iyer-au/retail-world-models}, 2026.\fi
\bibitem{smithwallis2009} J.~Smith and K.~F. Wallis. A simple explanation of the forecast combination puzzle. \emph{Oxford Bulletin of Economics and Statistics}, 71(3):331--355, 2009.
\bibitem{wang2023} X.~Wang, R.~J. Hyndman, F.~Li, and Y.~Kang. Forecast combinations: An over 50-year review. \emph{International Journal of Forecasting}, 39(4):1518--1547, 2023.
\bibitem{vovk2005} V.~Vovk, A.~Gammerman, and G.~Shafer. \emph{Algorithmic Learning in a Random World}. Springer, 2005.
\bibitem{favorita} Corporaci\'on Favorita. Corporaci\'on Favorita grocery sales forecasting. Kaggle competition, 2017. \url{https://www.kaggle.com/c/favorita-grocery-sales-forecasting}.
\bibitem{timestral} Mistral AI. Timestral 1.0: pretrained time-series forecasting models, 15 October 2026. Code \url{https://github.com/mistralai/timestral}, commit a534bdb; models \texttt{mistralai/Timestral-28M} and \texttt{mistralai/Timestral-209M} on Hugging Face.
\bibitem{hansen2011} P.~R. Hansen, A.~Lunde, and J.~M. Nason. The model confidence set. \emph{Econometrica}, 79(2):453--497, 2011.
\end{thebibliography}


\clearpage
\appendix
\begin{table*}[t]
\centering
\caption{\textbf{Every constraint the trained models and the ensemble work
within, with its value and where it lives in the code.}}
\label{tab:constraints}
\footnotesize
\renewcommand{\arraystretch}{1.15}
\begin{tabular}{p{0.29\textwidth}p{0.29\textwidth}p{0.34\textwidth}}
\toprule
\textbf{Constraint} & \textbf{Value} & \textbf{Where} \\
\midrule
Log forecast bound (world and lift model) & $[-10, 10]$ around the item's level & \nolinkurl{src/rwm/model/state_model.py} readout \\
Lift when nothing is planned & Exactly zero & Readout; tested in \texttt{tests/} \\
Baseline under a changed plan & Unchanged & Readout; tested in \texttt{tests/} \\
Regular price & Highest price in the last 12 periods & Plan features \\
Brand-switching share $\alpha$ & One per run, fitted on the last 8 training weeks & Dominick's world model \\
Forecasts & Not below zero & Every model \\
Ensemble weights & One third each, fixed 13 October 2026 & \texttt{scripts/revision\_analysis.py} \\
Interval pooling & Horizons pooled below 50 earlier errors per horizon & \texttt{src/rwm/evaluation/compare.py} \\
Interval ends & Rounded outwards to whole units (not Favorita) & \texttt{src/rwm/evaluation/compare.py} \\
TimesFM covariate fit & Groups of 2,048 series & \texttt{src/rwm/prior\_work/} \\
Timestral context and batch & 512 periods; batch 256, halved on running out of GPU memory & \nolinkurl{src/rwm/prior_work/timestral_zero_shot.py} \\
\bottomrule
\end{tabular}
\end{table*}

\section{Constraints the models work within}
\label{sec:constraints}



Table~\ref{tab:constraints} lists the constraints. The world model's settings are
in Table~\ref{tab:final-settings}.

\label{sec:supplement}
